%% ---------------------------------------------------------------------------
%% proposal.tex
%%
%% Research Proposal, main document.
%%
%% ---------------------------------------------------------------------------
\documentclass[12pt,letterpaper]{article}
\usepackage[english]{babel}     % supports english, but default is
% \usepackage[spanish]{babel}
% include this if you want to import graphics files with /includegraphics

\usepackage{longtable}
\usepackage{ifpdf}
\usepackage[table]{xcolor}

\usepackage{anysize}
\marginsize{2.5cm}{2.5cm}{1cm}{1cm}
\usepackage{textcomp}
\usepackage{url}
\bibliographystyle{unsrt}
\usepackage{graphics}
\usepackage{amssymb}
\usepackage{graphicx}
%\usepackage{slashbox}
\usepackage[latin1]{inputenc}
\usepackage{tikz}
\usetikzlibrary{arrows,positioning}
% Color and strikethrough



\usepackage{color}
\usepackage{soul}

\usepackage{array}
\usepackage{makecell}

\usepackage{sectsty}
\allsectionsfont{\sffamily}

\definecolor{dblue}{RGB}{0,102,153}
\newcommand{\dB}[1]{\textcolor{dblue}{\textbf{#1}}}



\setlength{\parskip}{1em}

% Nombre del Estudiante
\newcommand{\scriptAuthor}{Daniel Moya S�nchez}
 
% T��tulo de la tesis
\newcommand{\scriptTitle}{Progress report \#2 for the project: Design of (ASIPs) for Approximate Computing} 


% Keywords
\newcommand{\scriptKeywords}{key, words, ...}

% Para el PDF (cambiar si se desea otras cosas a lo indicado arriba
\newcommand{\pdfAuthor}{\scriptAuthor}
\newcommand{\pdfTitle}{\scriptTitle} 
\newcommand{\pdfKeywords}{\scriptKeywords}


\tikzset{
    mynode/.style={rectangle,rounded corners,draw=black, top color=white, bottom color=yellow!50,very thick, inner sep=1em, minimum size=3em, text centered},
    myarrow/.style={->, >=latex', shorten >=1pt, thick},
    mylabel/.style={text width=7em, text centered} 
} 

\begin{document}
 
 \graphicspath{{./}{./fig/}}

 %% ---------------------------------------------------------------------------
%% titlepage.tex
%%
%% Title page
%%
%% ---------------------------------------------------------------------------

\thispagestyle{empty} 

\begin{center}

\textsc{\LARGE Instituto Tecnol\'ogico de Costa Rica} \\
\textsc{\Large \'Area Acad\'emica de Ingenier\'ia en Computadores}

\textsc{\Large Proyecto de Dise\~no en Ingenier\'ia en Computadores}


\par\vspace{20mm}

\includegraphics[scale=0.24]{logoTEC}

\par\vspace*{\fill}

{\LARGE\bf{\textsf{ \Huge \scriptTitle}}}

\par\vspace*{\fill}

Chair for Embedded System (CES) \\
Kalrsruhe Institute of Technology (KIT) \\
Period: 19/02/2018 (week 2) - 23/02/2018 (week 4) 
%Master Thesis {\sf Proposal} \\ 
%in fulfillment of the requirements for the degree of
%Plan de Proyecto
%Master of Science in Electronics Engineering \\
%Emphasis on Embedded Systems

\par\vspace{20mm}

\textsc{\Large \scriptAuthor}

\vspace*{\fill}

{\today}

\end{center}
\newpage 
\cleardoublepage  


%  \tableofcontents

 \clearpage

\section{Performed activities} \label{act}

\begin{enumerate}
 \item \textbf{Familiarization with the software platform (ID 01)}:  More laboratory scripts have been
 studied to get to know the software tools from which the ASIPs are going to be built. However, the
 activity is still being affected by access rights and general configuration of the environment.
 The laboratory sessions are expected to be completed by the next report, since they are almost
 finished by now. 
  
 \item \textbf{Find appropiate error-tolerant applications (ID 02)}: Six possible applications have
 been found and analyzed for their approximate adaptation, this possibilities have been delivered
 to professor Jorge Castro, further discussion and his approval is being waited. 
 
 
 \item \textbf{Write Design document (ID 06)}: The final design document was send to the professor.

 \end{enumerate}
 

\section{Scope Changes} \label{change}

No scope changes have been made.

\section{Earned Value analysis}

Table \ref{tab:gain} summarizes the gained value analysis. 

\begin{table}[h!]
\begin{center}
\caption{Earned Value}
\resizebox{\textwidth}{!}{\begin{tabular}{|c|c|c|c|c|c|c|c|c|c|c|c|c|c|} 
 \hline
\makecell{Activity\\ID}	&Activity	&Budget	&\makecell{\%Planned\\Value}	&PV	&AC	&\makecell{\%Completed\\work}	&EV	
&CPI	&SPI	&\makecell{Initial planned\\date}	&\makecell{Ending\\date}	&\makecell{Initial real\\date}	&\makecell{Real ending} \\ \hline


01	&\makecell{Get to\\ know the\\ software\\ platform}	&32	&100\%	&32	&30	&95\%	&$30.4$
&$1.01$	&0.95	&Week 1	&Week 3	&Week 1	&- \\ \hline

02	&\makecell{Find appropiate\\ error-tolerant\\ applications}	&32	&100\%	&32	&15	&70\%	&22.4
&1.49	&0.7	&Week 4	&Week 6	&Week 4	&- \\ \hline

06	& \makecell{Write Design\\ document}	&8	&100\%	&8	&7	&100\%	&8
&1.14	&1	&\makecell{Week 3}	&\makecell{Week 4}	&\makecell{Week 3}	&\makecell{Week 5} \\ \hline

& Total	& &\%PVtotal	&PV total	&AC total 	&\%CW total	&EV total
&CPI total	&SPI total	&	&	&	& \\ \hline
\end{tabular}}
\label{tab:gain}
\end{center}
\end{table}




\section{Difficulties Encountered}

\begin{itemize}
 \item The server errors have slown down the progress of the activity 01, which should have been finished by the 
 last report.
 \item The personal computer used for work got a fan malfunctioning, which prevented the use of Linux (since Windows does
 have a fan control driver contrary to Linux) and contribuited to the slow progress on the laboratory sessions.
 \item Supervisor Jorge Castro was busy on week 6, situation that delayed the possible approximate applications review. 
\end{itemize}




\section{Hard Skills Required/Acquired}

Knowledge in the following software frameworks has been acquired:  Dlxsim, ModelSim and Xilinx ISE. 


\section{Soft Skills Required/Acquired}

The following soft skills have been excercised:

\begin{itemize}
 \item Communication: Weekly remote communication has been performed with Jorge Castro for the guidance of this project, and
 with Sajjad Hussain to request technical aid in the server. With both, swift communication was achieved, each topic that
 was talked was resolved or clarified in very few messages (one or two at most).
 
 \item Self-Motivation: Given that there is no direct round-the-clock supervision, self-motivation has been key in working 
 continously in the laboratory sessions and in the search of possible approximate applications.
 
\end{itemize}


\section{Lessons Learned}

It is very important to have tasks than can be parallelized because whenever something out of control (like the computer
malfunctioning) stops a task, one can still advance in another activity so the general impact on the schedule is not so big. 



% Referencias del Background y el Related Work
\bibliographystyle{sty/plainurl}
\bibliography{references}



\end{document}

